% ============ 第2章 §2.18 Carleson插值问题 (I)(原书 pp.179-185) ============
\section{Carleson 插值问题 (I)}\label{sec:2.18}
\index{一致分离}
\index{分离}


对 D 中相异的无限个点 $\{z_{j}\}_{j=1}^{\infty}$ 与 $\mathbb{C}$ 中无限个数据 $\{w_{j}\}_{j=1}^{\infty}$,求属于 $H^{p}$ 且满足
\begin{equation*}
f(z_{j}) = w_{j} \quad (1 \leq j < \infty)
\end{equation*}
的函数 f,与 $\{z_{j}\}$ 为有限个点的情况不同,不一定总存在.

如 2.17 节那样,对相异的 $\{z_{j}\}_{j=1}^{n}$ 定义 Blaschke 积 $B_{n}(z)$ 与 $B_{nj}(z)$,令
\begin{equation*}
f_{n}(z) = \sum_{j=1}^{n} w_{j}\frac{B_{nj}(z)}{B_{nj}(z_{j})}.
\end{equation*}
此时 $f_{n}(z_{j}) = w_{j}$($1 \leq j \leq n$).下面的命题~\ref{prop:2.18.1} 给出无限个点的插值问题的一般解答.

\begin{proposition}\label{prop:2.18.1}
设 $1 \leq p \leq \infty$ 且 $1/p + 1/q = 1$.又设 $\{z_{j}\}_{j=1}^{\infty} \subset D$ 且 $\{w_{j}\}_{j=1}^{\infty} \subset \mathbb{C}$.此时存在满足 $f(z_{j}) = w_{j}$($1 \leq j < \infty$)的 $f \in H^{p}$ 的充要条件是,存在常数 $\gamma > 0$,使得对所有 $g \in H^{q}$,
\begin{equation*}
\left| \sum_{j=1}^{n} \frac{w_{j}}{B_{nj}(z_{j})}(1 - |z_{j}|^{2}) g(z_{j}) \right| \leq \gamma\|g\|_{q} \quad (n = 1, 2, \dots).
\end{equation*}
\end{proposition}

\begin{proof}
设 $w = (w_{j})_{j=1}^{\infty}$,对 $n \geq 1$,令
\begin{equation*}
m_{p,n}(w) = \inf \{\|f_{n} + B_{n}h\|_{p}; h \in H^{p}\},
\end{equation*}
由定理~\ref{thm:2.17.1},
\begin{align*}
m_{p,n}(w) &= \sup \left\{\left| \sum_{j=1}^{n} \frac{w_{j}}{B_{nj}(z_{j})}(1 - |z_{j}|^{2}) g(z_{j}) \right|; \right. \\
& \left. g \in H^{q} \ \text{且} \ \|g\|_{q} \leq 1 \right\}.
\end{align*}
设 $\gamma = \sup_{n}m_{p,n}(w) < \infty$.由定理~\ref{thm:2.14.2},存在 $F_{n} \in f_{n} + B_{n}H^{p}$ 且 $\|F_{n}\|_{p} = m_{p,n}(w)$.此时 $\|F_{n}\|_{p} \leq \gamma$($n \geq 1$)且 $F_{n}(z_{j}) = w_{j}$($1 \leq j \leq n$).由定理~\ref{thm:2.13.2} 与定理~\ref{thm:2.13.3},$H^{p}$ 的单位球是 *弱紧的,故 $\{F_{n}\}$ 存在子列 $\{F_{n_{k}}\}$,$F_{n_{k}}$ 在 *弱拓扑下收敛到某个 $f \in H^{p}$.由命题 2.9.4,
\begin{equation*}
\phi_{j}(h) = h(z_{j}) = \frac{1}{2\pi i}\int_{|z|=1} \frac{h(z)}{z - z_{j}}\dd z \quad (h \in H^{p}),
\end{equation*}
因 $(z - z_{j})^{-1} \in L^{q} \cap C$,$\phi_{j}$ 是 *弱连续的.于是 $f(z_{j}) = w_{j}$($1 \leq j < \infty$).反之,若存在满足 $f(z_{j}) = w_{j}$($1 \leq j < \infty$)的 $f \in H^{p}$,则
\begin{equation*}
\inf \{\|f + B_{n}h\|_{p}; h \in H^{p}\} \leq \|f\|_{p},
\end{equation*}
故 $\sup_{n}m_{p,n}(w) < \infty$.定理得证.
\end{proof}

命题~\ref{prop:2.18.1} 作为无限个点的插值问题的解答完全不充分.下面的命题~\ref{prop:2.18.2} 给出对特殊数据 $\{w_{j}\}_{j=1}^{\infty}$ 的解答.设 $\delta_{\ell j} = 0$($\ell \neq j$)且 $\delta_{\ell j} = 1$($\ell = j$).

\begin{proposition}\label{prop:2.18.2}
设 $0 < p \leq \infty$ 且 $(w_{j})_{j=1}^{\infty} = (\delta_{\ell j})_{j=1}^{\infty}$.存在满足
\begin{equation*}
f(z_{j}) = w_{j} \quad (1 \leq j < \infty)
\end{equation*}
的 $f \in H^{p}$ 的充要条件是 $\sum_{j=1}^{\infty}(1 - |z_{j}|) < \infty$.
\end{proposition}

\begin{proof}
若 $\sum_{j=1}^{\infty}(1 - |z_{j}|) < \infty$,由定理~\ref{thm:2.6.1},存在对 $\{z_{j}\}_{j=1}^{\infty}$ 的 Blaschke 积 B.对任意自然数 $\ell$,令
\begin{equation*}
B_{\ell}(z) = B(z)\frac{z_{\ell}}{|z_{\ell}|}\frac{1 - \bar{z}_{\ell}z}{z_{\ell} - z},
\end{equation*}
令
\begin{equation*}
f(z) = B_{\ell}(z)/B_{\ell}(z_{\ell}),
\end{equation*}
则 $f \in H^{p}$ 且 $f(z_{j}) = \delta_{\ell j}$($1 \leq j < \infty$).反之,若存在满足 $f(z_{j}) = \delta_{\ell j}$($1 \leq j < \infty$)的 $f \in H^{p}$,由定理~\ref{thm:2.6.3} 得 $\sum_{j=1}^{\infty}(1 - |z_{j}|) < \infty$.
\end{proof}

令
\begin{equation*}
\ell^{0} = \{(w_{j}); (w_{j}) = (\delta_{\ell j}), 1 \leq \ell < \infty\},
\end{equation*}
若 $\sum_{j=1}^{\infty}(1 - |z_{j}|) < \infty$,由命题~\ref{prop:2.18.2},$\{(f(z_{j})); f \in H^{p}\} \supset \ell^{0}$.设 $\ell^{\infty} = \{(w_{j}); \sup_{j}|w_{j}| < \infty\}$,则
\begin{equation*}
\ell^{\infty} \supset \{(f(z_{j})); f \in H^{\infty}\} \supset \ell^{0}.
\end{equation*}
弄清 $\{(f(z_{j})); f \in H^{\infty}\}$ 是 $\ell^{\infty}$ 的怎样的子空间是最终目标,Carleson \cite{ch2:16} 给出了 $\{(f(z_{j})); f \in H^{\infty}\} = \ell^{\infty}$ 成立的充要条件.这一 Carleson 的结果被 Shapiro 与 Shields \cite{ch2:17} 推广到 $1 \leq p < \infty$,被 Kabaila \cite{ch2:18} 推广到 $0 < p < 1$.但他们研究的不是 $\{(f(z_{j})); f \in H^{p}\}$,而是 $\{((1 - |z_{j}|^{2})^{1/p}f(z_{j})); f \in H^{p}\}$.对 $0 < p < \infty$,设 $\ell^{p} = \{(w_{j}); \sum_{j=1}^{\infty}|w_{j}|^{p} < \infty\}$.

若 $\{z_{j}\}_{j=1}^{\infty} \subset D$ 满足 $\sum_{j=1}^{\infty}(1 - |z_{j}|) < \infty$,由定理~\ref{thm:2.6.1} 存在 Blaschke 积
\begin{equation*}
B(z) = \prod_{j=1}^{\infty} \frac{|z_{j}|}{z_{j}}\frac{z_{j} - z}{1 - \bar{z}_{j}z}.
\end{equation*}
\begin{equation*}
\begin{aligned}
\rho_{k,n} &= \prod_{\substack{j=1\\ j\neq k}}^{n}\left|\frac{z_{j} - z_{k}}{1 - \bar{z}_{j}z_{k}}\right|\\
\rho_{k} &= \prod_{\substack{j=1\\ j\neq k}}^{\infty}\left|\frac{z_{j} - z_{k}}{1 - \bar{z}_{j}z_{k}}\right|
\end{aligned} \quad (1 \leq k \leq n),
\end{equation*}
则 $\rho_{k,n} \geq \rho_{k,n+1}$ 且 $\lim_{n \to \infty}\rho_{k,n} = \rho_{k}$($k \geq 1$).一般 $\rho_{k} > 0$ 未必成立,但当 $\rho_{k} > 0$($k \geq 1$)时称 $\{z_{j}\}_{j=1}^{\infty}$ 是分离的.当 $\inf_{k}\rho_{k} > 0$ 时称 $\{z_{j}\}_{j=1}^{\infty}$ 是一致分离的.

\begin{theorem}\label{thm:2.18.3}
设 $0 < p \leq \infty$.$\{((1 - |z_{j}|^{2})^{1/p}f(z_{j})); f \in H^{p}\} = \ell^{p}$ 成立的充要条件是 $\{z_{j}\}_{j=1}^{\infty}$ 一致分离.
\end{theorem}

设 $\ell^{p}((1 - |z|^{2})^{1/p}) = \{(w_{j}); \sum_{j=1}^{\infty}(1 - |z_{j}|^{2})|w_{j}|^{p} < \infty\}$,则 $\{((1 - |z_{j}|^{2})^{1/p}f(z_{j})); f \in H^{p}\} = \ell^{p}$ 与 $\{(f(z_{j})); f \in H^{p}\} = \ell^{p}((1 - |z|^{2})^{1/p})$ 等价.

下面证明定理~\ref{thm:2.18.3},p = 2 的情形(准确地说充分条件证明的一部分)证明后,$p \neq 2$ 的情形可由之导出.下面的命题~\ref{prop:2.18.4} 的证明将在 2.19 节给出,本节利用它给出目标定理~\ref{thm:2.18.3} 的证明.

\begin{proposition}\label{prop:2.18.4}
$\{((1 - |z_{j}|^{2})^{1/2}f(z_{j})); f \in H^{2}\} = \ell^{2}$ 成立的充要条件是 $\{z_{j}\}_{j=1}^{\infty}$ 一致分离.
\end{proposition}

\begin{proof}
引理~\ref{lem:2.18.5} 证明必要性.若 $\inf_{j}\rho_{j} > 0$,则 $\{((1 - |z_{j}|^{2})^{1/2}f(z_{j})); f \in H^{2}\} = \ell^{2}$ 将在 2.19 节证明.
\end{proof}

\begin{lemma}\label{lem:2.18.5}
设 $0 < p \leq \infty$.若 $\{((1 - |z_{j}|^{2})^{1/p}f(z_{j})); f \in H^{p}\} = \ell^{p}$,则 $\inf_{j}\rho_{j} > 0$.
\end{lemma}

\begin{proof}
令 $T_{p}(f) = ((1 - |z_{j}|^{2})^{1/p}f(z_{j}))$,由假设 $T_{p}(H^{p}) = \ell^{p}$.设 $f_{n} \in H^{p}$,$f_{n} \to f$ 且
\begin{equation*}
T_{p}(f_{n}) = ((1 - |z_{j}|^{2})^{1/p}f_{n}(z_{j})) \to w \in \ell^{p},
\end{equation*}
由 $f_{n}(z_{j}) \to f(z_{j})$($j \geq 1$),$T_{p}(f) = w$.这表明 $T_{p}$ 是闭算子,由闭图像定理,$T_{p}$ 是从 $H^{p}$ 到 $\ell^{p}$ 的有界算子.

令
\begin{equation*}
\ker T_{p} = \{f \in H^{p}; f(z_{j}) = 0 \ (j \geq 1)\},
\end{equation*}
因 $\ker T_{p}$ 是算子 $T_{p}$ 的核,由 $T_{p}(H^{p}) = \ell^{p}$,可从 $T_{p}$ 定义从 $H^{p}/\ker T_{p}$ 到 $\ell^{p}$ 的双射算子 $\underline{T_{p}}$,即令 $\underline{T_{p}}(f + \ker T_{p}) = T_{p}(f)$($f \in H^{p}$).此时 $\underline{T_{p}}$ 可逆,故对任意 $w \in \ell^{p}$,
\begin{equation*}
(1 - |z_{j}|^{2})^{1/p}f(z_{j}) = w_{j} \quad (j \geq 1)
\end{equation*}
且设 $\gamma = \|\underline{T_{p}}^{-1}\|$,则存在满足
\begin{equation*}
\|f\|_{p} \leq \gamma\|w\|_{p}
\end{equation*}
的 $f \in H^{p}$.于是对每个 k,存在 $f_{k} \in H^{p}$,使得 $\|f_{k}\|_{p} \leq \gamma$ 且
\begin{equation*}
f_{k}(z_{j}) = \left\{ \begin{array}{l l}
(1 - |z_{k}|^{2})^{-1/p} & (j = k) \\
0 & (j \neq k)
\end{array} \right..
\end{equation*}
对 n > k,令
\begin{equation*}
F_{nk}(z) = f_{k}(z)\prod_{\substack{j=1\\ j\neq k}}^{n}\frac{1 - \bar{z}_{j}z}{z_{j} - z},
\end{equation*}
则 $F_{nk} \in H^{p}$ 且 $\|F_{nk}\|_{p} = \|f_{k}\|_{p} \leq \gamma$.由引理~\ref{lem:2.10.1},
\begin{equation*}
|F_{nk}(z_{k})| \leq 2^{1/p}\|F_{nk}\|_{p}(1 - |z_{k}|)^{-1/p},
\end{equation*}
故
\begin{equation*}
\left|\prod_{\substack{j=1\\ j\neq k}}^{n}\frac{1 - \bar{z}_{j}z_{k}}{z_{j} - z_{k}}\right| = |f_{k}(z_{k})|^{-1}|F_{nk}(z_{k})| \leq 2^{1/p}\gamma.
\end{equation*}
这证明了 $\inf_{j}\rho_{j} > 0$.
\end{proof}

\begin{lemma}\label{lem:2.18.6}
当 $0 < p \leq \infty$ 时,若 $\inf_{j}\rho_{j} > 0$,则 $\{((1 - |z_{j}|^{2})^{1/p}f(z_{j})); f \in H^{p}\} \subset \ell^{p}$.
\end{lemma}

\begin{proof}
p = $\infty$ 时显然,可设 0 < p < $\infty$.由定理~\ref{thm:2.6.4},$f \in H^{p}$ 可写成 $f(z) = B(z)[F(z)]^{2/p}$,其中 B 是 Blaschke 积,F 是 D 中无零点的 $H^{2}$ 函数.于是由命题~\ref{prop:2.18.4},
\begin{equation*}
\sum_{j=1}^{\infty}(1 - |z_{j}|^{2})|f(z_{j})|^{p} \leq \sum_{j=1}^{\infty}(1 - |z_{j}|^{2})|F(z_{j})|^{2} < \infty.
\end{equation*}
\end{proof}

\begin{lemma}\label{lem:2.18.7}
当 $1 \leq p \leq \infty$ 时,若 $\inf_{j}\rho_{j} > 0$,则 $\{((1 - |z_{j}|^{2})^{1/p}f(z_{j})); f \in H^{p}\} \supset \ell^{p}$.
\end{lemma}

\begin{proof}
由定理~\ref{thm:2.17.1},存在满足
\begin{equation*}
g_{n}(z_{j}) = (1 - |z_{j}|^{2})^{-1/p}w_{j} \quad (1 \leq j \leq n)
\end{equation*}
的最小范数的 $g_{n} \in H^{p}$.又当 $p \neq 1$ 时,设 $1/p + 1/q = 1$,则存在满足 $\|f\|_{q} = 1$ 的 $f \in H^{q}$,使得
\begin{equation*}
\|g_{n}\|_{p} = \left| \sum_{j=1}^{n} \frac{w_{j}(1 - |z_{j}|^{2})^{1/q}}{B_{nj}(z_{j})} f(z_{j}) \right|.
\end{equation*}
将
\begin{equation*}
\frac{|w_{j}|(1 - |z_{j}|^{2})^{1/q}}{|B_{nj}(z_{j})|}|f(z_{j})| \leq \rho_{j}^{-1}\cdot|w_{j}|\cdot(1 - |z_{j}|^{2})^{1/q}|f(z_{j})|
\end{equation*}
代入,由 Hölder 不等式,
\begin{align*}
\|g_{n}\|_{p} &\leq (\inf_{j}\rho_{j})^{-1}\|w\|_{p}\left\{\sum_{j=1}^{n}(1 - |z_{j}|^{2})|f(z_{j})|^{q}\right\}^{1/q} \\
&\leq \gamma(\inf_{j}\rho_{j})^{-1}\|w\|_{p}\|f\|_{q} \\
&\leq \gamma(\inf_{j}\rho_{j})^{-1}\|w\|_{p}.
\end{align*}
这里 $\gamma$ 的存在由引理~\ref{lem:2.18.6} 与闭图像定理保证.

由引理~\ref{lem:2.10.1},$\{g_{n}\}$ 是正规族,故存在 $g \in H(D)$,对每个 R < 1,$\{g_{n}\}$ 的某个子列在 $\{z; |z| < R\}$ 上一致收敛.
\end{proof}

\begin{lemma}\label{lem:2.18.8}
当 $0 < p \leq 1$ 时,若 $\inf_{j}\rho_{j} > 0$,则 $\{((1 - |z_{j}|^{2})^{1/p}f(z_{j})); f \in H^{p}\} \supset \ell^{p}$.
\end{lemma}

\begin{proof}
令 $f_{j}(z) = (1 - \bar{z}_{j}z)^{-2/p}$,
\begin{equation*}
f(z) = \sum_{j=1}^{\infty}(1 - |z_{j}|^{2})^{1/p}w_{j}[B_{j}(z_{j})]^{-1}B_{j}(z)f_{j}(z).
\end{equation*}
设 $(w_{j}) \in \ell^{p}$ 且 $\inf_{j}\rho_{j} > 0$,则对任意 r,f 在 $|z| \leq r < 1$ 上一致收敛,故 $f \in H(D)$.进而由 Minkowski 不等式,
\begin{align*}
M_{p}^{p}(r, f) &\leq (\inf_{j}\rho_{j})^{-p}\sum_{j=1}^{\infty}(1 - |z_{j}|^{2})|w_{j}|^{p}M_{p}^{p}(r, f_{j}) \\
&\leq (\inf_{j}\rho_{j})^{-p}\|w\|_{p}^{p}.
\end{align*}
\end{proof}

定理 2.18.2 的证明 必要条件是引理~\ref{lem:2.18.5}.证明充分性.设 $\inf_{j}\rho_{j} > 0$,由利用命题~\ref{prop:2.18.4}(尚未证明的充分性)证明的引理~\ref{lem:2.18.6},当 0 < p $\leq \infty$ 时
\begin{equation*}
\{((1 - |z_{j}|^{2})^{1/p}f(z_{j})); f \in H^{p}\} \subset \ell^{p}.
\end{equation*}
关于相反的包含关系,当 0 < p $\leq$ 1 时,由引理~\ref{lem:2.18.8},
\begin{equation*}
\{((1 - |z_{j}|^{2})^{1/p}f(z_{j})); f \in H^{p}\} \supset \ell^{p}.
\end{equation*}
当 $1 \leq p \leq \infty$ 时由引理~\ref{lem:2.18.7} 证明.定理 2.18.2 的证明中还剩下的只有命题~\ref{prop:2.18.4} 的充分性.
